\documentclass[10.5pt]{article}
\usepackage[margin=1in]{geometry}
\usepackage{parskip}
\usepackage{xcolor}
\usepackage{hyperref}
\usepackage{fontspec} % remove/replace with [T1]{fontenc} if not using XeLaTeX/LuaLaTeX

\usepackage{xcolor}
\usepackage{marginnote}
\newcommand{\TODO}[1]{%
  \marginnote{\color{red}\textbf{TODO}}% Put 'TODO' in margin
  {\color{red}#1}% Put the associated text in red in the body
}

% All editable fields (recipient, position, blanks, etc.) live in macros.tex
\input{macros}

\setlength{\parindent}{0pt}
\pagestyle{empty}

\begin{document}

% ---------- Header ----------
\begin{flushright}
\textbf{\ApplicantName}\\ 
% \ApplicantTitle \\
\ApplicantInstitutionNameOne\\
\ApplicantInstitutionNameTwo\\
\ApplicantInstitutionAddressLine\\ \ApplicantInstitutionCityLine\\
\end{flushright}

\vspace{0.25em}
\RecipientNameOne \hfill \LetterDate \\
\RecipientInstitution \\
% \RecipientDept\\
% \RecipientTitle\\
% \RecipientAddress

% \vspace{1em}
Dear \RecipientNameTwo,

\vspace{0.5em}

I am writing to express my strong interest in the \PositionTitle\ at \RecipientInstitution. 
I am currently a \ApplicantTitle at the \ApplicantInstitutionNameTwo.
During my PhD, my advisor moved from the University of Michigan-Ann Arbor to the University of Zürich. This change propelled me to develop a unique, independent research program focused on constraining cosmology using electromagnetic and gravitational wave data. 
In the development of this program, I have developed extensive expertise in astronomical instrumentation, gravitational-wave cosmology, and multi-messenger astrophysics, all of which align closely with \HostGroupDescriptor. I believe that my work, along with my experience and continued engagement in outreach and science education will positively contribute to \RecipientInstitution.

% \textcolor{red}{[Something about the move to Zurich, that was actually to Chile]}
Shortly after beginning my studies in Zürich, I moved to La Serena, Chile to support the Vera C. Rubin Observatory through the final stages of construction and early science operations.
While in Chile, I led the commissioning of the LSST-Camera focal plane array, including routine operation, intervention, and optimization of the LSST-Camera CCDs. 
% \textcolor{red}{Something that connects to the host institution.}
\textcolor{\tcolor}{I would like to transfer this expertise to Fermilab - leading image sensors characterization for future experiments, contributing to both smaller telescopes like SkipperCam, and pathfinding research for larger facilities such as the Spectroscopic Stage-V Experiment.}

% testing robotic positioners for the DESI instrument and future Stage-V spectroscopic surveys, and in a previous role, delivering the fiber run and spectrograph cameras for the LLAMAS spectrograph for the Magellan Baade telescope. 

In parallel to commissioning the LSST-Camera, I led several distinct observational programs in multi-messenger astrophysics, with a focus on constraining cosmological parameters using observations from optical telescopes and gravitational wave detectors. This includes leading Target-of-Opportunity searches with Vera C. Rubin Observatory, the Dark Energy Camera, and multi-instrument searches with the Dark Energy Camera and the Prime Focus Spectrograph. 
% \textcolor{red}{Something here that connects to the host institution.} 
\textcolor{\tcolor}{Bringing this expertise to Fermilab would help to sustain the aforementioned observational programs, while ensuring that forthcoming facilities led by Fermilab, including SkipperCam and the Spectroscopic Stage-V Experiment, are informed by my multi-messenger expertise.} 
%, and the LIGO-Virgo-Kagra gravitational wave observatory. 
% In addition to maximizing data from existing facilities, I am preparing analysis techniques and observing strategies for future observatories, focusing on Stage-V spectroscopic facilities, proposed spectroscopic facilities, and Rubin Observatory to detect the electromagnetic signatures, and forecasted gravitational wave data from Cosmic Explorer, the Einstein Telescope, and the Laser Interferometer Space Antenna. 

% Through these activities, I have become an active member of the leading ground based electromagnetic observatories 
% (Legacy Survey of Space and Time Dark Energy Science Collaboration, Rubin Observatory Project) 
% and gravitational wave observatories
% (LIGO-Virgo-Kagra, the Laser Interferometer Space Antenna).
% , and have established a group of collaborators spanning the United States, Chile, France, Italy, Japan, Great Britain, Germany, Switzerland, and Spain.

Throughout my research career, I have prioritized public engagement and outreach as a core complement to my scientific work. 
Recently, I led an exhibition at the University of Zürich's Science pavilion during the City of Zürich's long night of museums. Leading `Ask a Scientist' sessions, focused on observational cosmology and gravitational waves, reinforced that translating complex physics concepts for non-specialist audiences not only broadens the impact of my research, but also sharpens my own understanding.
\textcolor{\tcolor}{I am committed to working with the Fermilab’s Education and Public Outreach program to share this enthusiasm with the Saturday Morning Physics program and other programs in the wider Chicagoland area.}

I have enclosed my \supplementList\ for your consideration. 
% Additionally, \ReferenceOne, \ReferenceTwo, and \ReferenceThree\ would be happy to provide references on my behalf. 
I welcome the opportunity to discuss how my expertise could contribute to the scientific mission of \RecipientInstitution\ and \HostGroupDescriptor.

\vspace{1.5em}
All the best,\\ \\ \\
\ApplicantName

\end{document}